\subsection*{Structured concurrent regions with \texttt{inthreads.\allowbreak{}async}}

The form \co@inthreads.async@ assigns different thread groups to separately progressing regions. Its predicate identifies the group responsible for each region. This is useful for a producer and consumer whose roles are fixed even though their progress depends on which buffers are ready.

\begin{lstlisting}[language=croqtile,caption={Structured concurrent regions with inthreads.async}]
parallel p1 by 2 : thread {
  inthreads.async (p1 == 0) {
  }
  inthreads.async (p1 == 1) {
  }
}
\end{lstlisting}

The distinction from \texttt{if} is structural, not just performance:

\begin{table}[htbp]
\centering\small
\begin{tabularx}{\linewidth}{@{}lXX@{}}
\toprule
 & \textbf{\texttt{if} (predicated execution)} & \textbf{\texttt{inthreads.\allowbreak{}async} (structured concurrency)} \\
\midrule
\textbf{Resolution} & Runtime --- every thread evaluates the condition & Compile time --- thread assignment is fixed \\
\textbf{Instruction streams} & One program; divergent threads masked & Separate programs per region \\
\textbf{PL analogy} & \texttt{if}/\texttt{else} in any language & \texttt{async}/\texttt{spawn} in structured concurrency (Trio, Go goroutines, Cilk) \\
\bottomrule
\end{tabularx}
\end{table}

The regions are lexically scoped, so the compiler can identify their participants from the program structure. The code does not create an unbounded number of tasks at runtime. Each region operates within the declared parallel scope.

Without \co@.async@, the region rejoins the surrounding execution before later code proceeds. Adding \co@.async@ allows other regions to advance independently. Any data handoff between them then needs an explicit synchronization protocol.

The figure below shows the effect. --- The producer's DMA and the consumer's MMA overlap in time:

Figure~\ref{fig:f07-fig1_role_comparison} compares the two schedules.

\begin{figure}[htbp]\centering
\includegraphics[width=\linewidth,height=0.56\textheight,keepaspectratio]{f07-fig1_role_comparison_light.png}
\caption{Four load/compute pairs are shown in both schedules. Concurrent producer and consumer regions overlap work while preserving dependencies.}
\label{fig:f07-fig1_role_comparison}
\end{figure}
